\section{A Brief History Of Economics}

Economics in its basic form began during the Bronze Age (4000-2500 BCE) with written documents in four areas of the
world: Sumer and Babylonia (3500-2500 BCE), the Indus River Valley Civilization (3300-1030 BCE), along the Yangtze
River in China, and Egypt's Nile Valley, beginning around 3500 BCE. Societies in these areas developed notation
systems using markings on clay tablets, papyrus,and other materials to account for crops, livestock, and land. These
accounting systems, arising in tandem with written language, eventually included methods for tracking property
transfers, recording debts and interest payments, calculating compound interest, and other economic tools still used
today. \v

From the third millennium BCE onward, Egyptian scribes recorded the collection of and redistribution of land and
goods. Sumerian traders developed methods to calculate compound interest over a period of months and years. The Code
of Hammurabi (circa 1810–1750 BCE), the earliest work of economic synthesis, specifies norms for economic activity
and provides a detailed framework for commerce, including business ethics for merchants and tradespeople. \v

The first millennium BCE saw the emergence of more detailed written treatises on economic thought and practice. The
Greek philosopher and poet Hesiod, writing in the eighth century BCE, laid out precepts for managing a farm in his
Works and Days. Athenian military leader, philosopher, and historian Xenophon built on this in Oikonomikon, a
treatise on the economic management of an estate. In Politics, Aristotle (circa 350 BCE) took these ideas further
still, concluding that while private ownership of property was preferred, the accumulation of wealth for its own sake
was ``dishonorable.'' \v

The Guanzi essays from China (circa the fourth century BCE) laid out one of the first explanations of supply and
demand pricing; the crucial roles of a well-managed money supply and a stable currency. Among key insights was the
notion that it was money, not armies, that ultimately won wars. \v

In Western Europe during the Middle Ages, economic theory was often blended with ethics, as seen in the work of
Thomas Aquinas (1225-1274) and others. Few of those writers went into the amount of detail that Ibn Khaldun
(1332-1406), Tunisian historian and philosopher, did. In Al-Muqaddimah, Ibn Khaldun analyzes economic issues such as
the perils of monopolies, the benefits of division of labor and the profit motive, and the rise and fall of economic
empires. The importance of his work was recognized by Machiavelli and Hegel, and many of his ideas prefigure those of
Adam Smith. \v

Today, Scottish thinker Adam Smith is widely credited with creating the field of modern economics. However, Smith was
inspired by French writers publishing in the mid-18th century, who shared his hatred of mercantilism. In fact, the
first methodical study of how economies work was undertaken by the French physiocrats, notably Quesnay and Mirabeau.
Smith took many of their ideas and expanded them into a thesis about how economies should work, as opposed to how
they do work. \v

Smith believed that competition was self-regulating and governments should take no part in business through tariffs,
taxes, or other means unless it was to protect free-market competition. Many economic theories today are, at least in
part, a reaction to Smith's pivotal work in the field, namely his 1776 masterpiece ``The Wealth of Nations''. In this
treatise, Smith laid out several of the mechanisms of capitalist production, free markets, and value. Smith showed
that individuals acting in their own self-interest could, as if guided by an ``invisible hand'', create social and
economic stability for all. \v

Even devout followers of Smith's ideas recognize that some of his theories were either flawed or have not aged well.
Smith distinguishes between ``productive labor'', such as manufacturing products that can be accumulated, and
``unproductive labor'', such as tasks performed by a ``menial servant'', the value of which ``perish[es] in the very
instant of their performance''. One could argue that in today's service-dominant economy, the excellent execution of
services creates value by strengthening a brand through goodwill and in numerous other ways. His assertion that
``equal quantities of labour, at all times and places, may be said to be of equal value to the labourer'' ignores the
psychological cost of working in hostile or exploitative environments. As an extension of this, Smith's labor theory
of value has also largely been abandoned. \v

Thomas Malthus and Karl Marx had decidedly poor reactions to Smith's treatise. Malthus was one of a group of economic
thinkers of the late 18th and early 19th centuries who were grappling with the challenges of emergent capitalism
following the French Revolution and the rising demands of a burgeoning middle class. Among his peers were three of
the greatest economic thinkers of the age, Jean-Baptiste Say, David Ricardo, and John Stuart Mill. \v

Malthus predicted that growing populations would outstrip the food supply. He was proved wrong, however, because he
didn't foresee technological innovations that would allow production to keep pace with a growing population.
Nonetheless, his work shifted the focus of economics to the scarcity of goods, rather than the demand for them. \v

This increased focus on scarcity led Marx to declare that the means of production were the most important components
of any economy. Marx took his ideas further and became convinced a class war was going to be sparked by the inherent
instabilities he saw in capitalism. However, Marx underestimated the flexibility of capitalism. Instead of creating a
clear division between two classes—owners and workers—the market economy created a mixed class wherein owners and
workers held the interests of both parties. Despite his overly rigid theory, Marx accurately predicted one trend:
businesses grow larger and more powerful to the degree that free-market capitalism allows. \v

As the ideas of wealth and scarcity developed in economics, economists turned their attention to more specific
questions about how markets operate and how market prices are determined. English economist William Stanley Jevons
(1835-1882), Austrian economist Carl Menger (1840-1921), and French economist Léon Walras (1834-1910) independently
developed a new perspective in economics known as marginalism. \v

Their key insight was that in practice, people aren't actually faced with big-picture decisions over entire general
classes of economic goods. Instead, they make their decisions around specific units of an economic good as they
choose to buy, sell, or produce each additional (or marginal) unit. In doing so, people balance the scarcity of each
good against the value of the use of the good at the margin. \v

These decisions explain, for example, why the price of an individual diamond is relatively higher than the price of
an individual unit of water. Though water is a basic need to live, it is often plentiful, and though diamonds are
often purely decorative, they are scarce. Marginalism quickly became, and remains, a central concept in economics. \v

Walras went on to mathematize his theory of marginal analysis and made models and theories that reflected what he
found. General equilibrium theory came from his work, as did the practice of expressing economic concepts
statistically and mathematically instead of just in prose. Alfred Marshall took the mathematical modeling of
economies to new heights, introducing many concepts that are still not widely understood, such as economies of scale,
marginal utility, and the real-cost paradigm. \v

It is nearly impossible to expose an economy to experimental rigor; therefore, economics is on the edge of science.
Through mathematical modeling, however, some economic theory has been rendered testable. The theories developed by
Walras, Marshall, and their successors would develop in the 20th century into the neoclassical school of
economics—defined by mathematical modeling and assumptions of rational actors and efficient markets. Later,
statistical methods were applied to economic data in the form of econometrics, giving economists the ability to
propose and test hypotheses empirically and in a methodologically rigorous manner. \v

John Maynard Keynes developed a new branch of economics known as Keynesian economics, or more generally as
macroeconomics. Keynes styled the economists who had come before him as ``classical'' economists, and he believed that
while their theories might apply to individual choices and goods markets, they did not adequately describe the
operation of the economy as a whole. \v

Instead of marginal units or even specific goods markets and prices, Keynesian macroeconomics presents the economy in
terms of large-scale aggregates that represent the rate of unemployment, aggregate demand, or average price-level
inflation for all goods. Keynes's theory says that governments can be powerful players in the economy and save it
from recession by implementing expansionary fiscal and monetary policy—manipulating government spending, taxing, and
money creation—in order to manage the economy. \v

By the mid-20th century, these two strands of thought—mathematical, marginalist microeconomics and Keynesian
macroeconomics—would rise to near-complete dominance of the field of economics throughout the Western world. This
became known as the neoclassical synthesis, which has since represented the mainstream of economic thought as taught
in universities and practiced by researchers and policymakers, with other perspectives labeled as heterodox economics. \v

Within the neoclassical synthesis, various streams of economic thought have developed, sometimes in opposition to one
another. The inherent tension between neoclassical microeconomics—which portrays free markets as efficient and
beneficial—and Keynesian macroeconomics—which views markets as inherently prone to calamitous failure—has led to
persistent academic and public policy disagreements, with different theories ascendant at different times. \v

Various economists and schools of thought have sought to refine, reinterpret, redact, and redefine both neoclassic
microeconomics and Keynesian macroeconomics. Most prominent is monetarism and the Chicago School, developed by Milton
Friedman, which retains neoclassical microeconomics and the Keynesian macroeconomic framework but shifts the emphasis
of macroeconomics from fiscal policy (favored by Keynes) to monetary policy. Monetarism was widely espoused through
the 80s to 00s. \v

Several different streams of economic theory and research have been proposed to resolve the tension between micro-
and macroeconomics by incorporating aspects or assumptions from microeconomics (such as rational expectations) into
macroeconomics or by further developing microeconomics to provide micro-foundations (such as price stickiness or
psychological factors) for Keynesian macroeconomics. In recent decades, this has led to the development of new
theories, such as behavioral economics, and to renewed interest in heterodox theories, such as Austrian-school
economics, which were previously relegated to the economic backwaters. \v

Classical economic theory and theory of markets, from Smith through Friedman, have rested largely on the assumption
that consumers are rational actors who behave in their own best interests. Current economists such as Richard Thaler,
Daniel Kahneman, Gary Becker, and the late Amos Tversky, have shown that people often do not act in their own best
material interests but allow themselves to be swayed by non-material, psychological factors, and biases. Behavioral
economics has helped to popularize a number of new heuristics and biases that make economic modeling and forecasting
more difficult than ever. \v

A rising cohort of economists has emphasized the importance of factoring in inequalities in income distribution and
social well-being when measuring the success of a given economic policy. Pre-eminent among them is Anthony Atkinson
(1944-2017), who focused on income redistribution within a given country, and Amartya Sen, professor of economics and
philosophy at Harvard University, whose work on global inequality won him the Nobel Prize for Economics in 1998.
Sen's work is also notable for reintroducing ethical behavior into his analysis. This concern ties Sen's thinking
back to the writing of the earliest economic thinkers, who saw over-accumulation of wealth by individuals or groups
as ultimately harmful to society. \v

Economic theory grew out of societies' need to account for resources, plan for the future, and exchange and allocate
goods. Over time, these basic accounting tools grew into financial models of increasing complexity, blending the
mathematics required to calculate compound interest, with ethics and moral philosophy. Economics as a system to
understand and control the material world and mitigate risk emerged and evolved across the globe in a staggered
fashion—the Fertile Crescent and Egypt, China and India, ancient Greece and the Arab world. \v

As societies grew wealthier and trade grew more complex, economic theory turned to the mathematics, statistics, and
computational modeling that economists use to help guide policymakers. The business cycle, booms and busts,
anti-inflation measures, and mortgage interest rates are outgrowths of economics. Understanding them helps the market
and government adjust for these variables. Balancing out the mathematical modeling approach is the study of factors
that are more difficult to quantify but crucial to understand—most notably, the foibles and unpredictability of human
psychology.

\section{Basic Terminology}

The word economy comes from the Greek word oikonomos, which means ``one who manages a household.'' At first, this
origin might seem peculiar. But in fact, households and economies have much in common. A household faces many
decisions. It must decide which household members do which tasks and what each member receives in return. In short, a
household must allocate its scarce resources among its various members, taking into account each member's abilities,
efforts, and desires. \v

Like a household, a society faces many decisions. It must find some way to decide what jobs will be done and who will
do them. It needs some people to grow food, other people to make clothing, and still others to design computer
software. Once society has allocated people (as well as land, buildings, and machines) to various jobs, it must also
allocate the goods and services they produce. It must decide who will eat caviar and who will eat potatoes. It must
decide who will drive a Ferrari and who will take the bus.

\bd[Goods \& Services]
\textbf{Goods} are tangible (or intangible) items that satisfy human wants and provide utility. A common distinction
is made between goods which are transferable, and \textbf{services}, which are not transferable.
\ed

\bd[Resource]
A \textbf{resource} is defined as anything used to produce goods and services that meet human needs and wants.
\ed

When thinking about the various goods in the economy, it is useful to group them by two characteristics:
``excludability'' which refers to if people can be prevented from using a good and ``rivalry'' which refers to if one
person's use of a unit of a good reduces another person's ability to use it.

\bd[Excludability]
\textbf{Excludability} is defined as the degree to which a good, service or resource can be limited to only paying
customers, or conversely, the degree to which a supplier, producer or other managing body can prevent ``free''
consumption of a good.
\ed

\bd[Excludable Good]
When a good caries the property of excludability is called an \textbf{excludable good}.
\ed

\bd[Rivalry]
\textbf{Rivalry} is defined as the degree to which the usage of a good by one person diminishes other people's use.
\ed

\bd[Rival Good]
When a good caries the property of rivalry is called a \textbf{rival good}.
\ed

Based on these two characteristics we can define four categories of goods: private, public, common, and club goods.

\bd[Private Good]
A \textbf{private good} is a good that is both excludable and rival.
\ed

Most goods in the economy are private goods. You don't get one unless you pay for it, and once you have it, you are
the only person who benefits.

\be
An example of a private good is an ice-cream. It is excludable because it is possible to prevent someone from eating
one (you just don't give it to her), and it is rival because if one person eats an ice-cream cone, another person
cannot eat the same cone.
\ee

\bd[Public Good]
A \textbf{public good} is a good that is neither excludable nor rival.
\ed

\be
An example of a public good is a tornado siren. Once the siren sounds, it is impossible to prevent any single person
from hearing it, so it is not excludable. Moreover, when one person gets the benefit of the warning, she does not
reduce the benefit to anyone else, so it is not rival.
\ee

\bd[Common Good]
A \textbf{common good} is a good that is not excludable but rival.
\ed

\be
An example of a common good is fishes in the oceans. They are rival because when one person catches some fish, fewer
fish are left for the next person to catch. But these fish are not an excludable good because it is hard to stop
fishermen from taking fish out of a vast ocean.
\ee

\bd[Club Good]
A \textbf{club good} is a good that is excludable but not rival.
\ed

\be
An example of a club good is satellite TV. If you don't pay the company offering the service, it can prevent you from
using it, making the good excludable. But your accessing the satellite signal does not diminish anyone else's ability
to access it, so the good is not rival.
\ee

The management of society's resources is important because resources are scarce, which means that society has limited
resources hence, it cannot produce all the goods and services people wish to have.

\bd[Scarcity]
\textbf{Scarcity} refers to a basic economics problem, the gap between limited resources and theoretically limitless
wants.
\ed

Scarcity as an economic concept refers to the basic fact of life that there exists only a finite amount of human and
nonhuman resources which the best technical knowledge is capable of using to produce only limited maximum amounts of
each economic good. This situation requires people to make decisions about how to allocate resources efficiently, in
order to satisfy basic needs and as many additional wants as possible. \v

Having introduced the concept of ``scarcity'', we can combine it with the definition of ``goods \& services'' in
order to define the so called ``economic good''.

\bd[Economic Good]
A good is an \textbf{economic good} if it is useful to people but scarce in relation to its demand so that human effort
is required to obtain it.
\ed

From the economic good we can define the concept of a commodity, which will be useful in later chapters.

\bd[Commodity]
A \textbf{commodity} is an economic good, usually a resource, that has full or substantial fungibility: that is, the
market treats instances of the good as equivalent or nearly so with no regard to who produced them.
\ed

Each economic good carries an economic value.

\bd[Economic Value]\label{eqn:economic_value}
\textbf{Economic value} is the value that person places on an economic good based on the benefit that they derive
from the good. It is often estimated based on the person's willingness to pay for the good, typically measured in
units of currency.
\ed

If the conditions of scarcity didn't exist and an infinite amount of every good could be produced or human wants
fully satisfied, there would be no economic goods. In economics, this situation of abundance is called ``post
scarcity'', although it does not mean that scarcity has been eliminated for all goods and services, but that all
people can easily have their basic survival needs met along with some significant proportion of their desires for
goods and services.

\bd[Post Scarcity]
\textbf{Post scarcity} is a theoretical economic situation in which most goods can be produced in great abundance
with minimal human labor needed, so that they become available to all very cheaply or even freely.
\ed

Economics is the study of how society manages its scarce resources.

\bd[Economics]
\textbf{Economics} is the science that studies the production, distribution, and consumption of economic goods,
services and commodities, and how individuals, businesses, governments, and nations make choices about how to
allocate resources.
\ed

In most societies, resources are allocated not by an all-powerful dictator but through the combined choices of
millions of households and firms. Economists therefore study how people make decisions: how much they work, what they
buy, how much they save, and how they invest their savings. Economists also study how people interact with one
another. For instance, they examine how the many buyers and sellers of a good together determine the price at which
the good is sold and the quantity that is sold. Finally, economists analyze the forces and trends that affect the
economy as a whole, including the growth in average income, the fraction of the population that cannot find work, and
the rate at which prices are rising. Economists try to address their subject with a scientist's objectivity. They
devise theories, collect data, and then analyze these data to verify or refute their theories.

\be
A theory might assert that high inflation arises when the government prints too much money. To test this theory, the
economist could collect and analyze data on prices and money from many different countries. If growth in the quantity
of money were unrelated to the rate of price increase, the economist would start to doubt the validity of this theory
of inflation. If money growth and inflation were correlated in international data, as in fact they are, the economist
would become more confident in the theory.
\ee

Although economists use theory and observation like other scientists, they face an obstacle that makes their task
especially challenging: they are not allowed to manipulate a system under investigation in order to generate useful
data but they usually have to make do with whatever data the world gives them. Economists pay close attention to the
natural experiments offered by history episodes. Studying these episodes is valuable because they give us insight
into the economy of the past and allow us to illustrate and evaluate economic theories of the present. \v

Economists, like other scientists, make assumptions. Assumptions can simplify the complex world and make it easier to
understand. The art in scientific thinking is deciding which assumptions to make. Similarly, economists use different
assumptions to answer different questions. Finally, economists use mathematical models, built with assumptions, to
learn about the world. \v

Economics can generally be broken down into the field of ``microeconomics'' which focuses on individual people and
businesses and in the field of ``macroeconomics'', which concentrates on the behavior of the economy as a whole.

\bd[Microeconomics]
\textbf{Microeconomics} is a branch of economics that studies the behavior of individuals and firms in making
decisions regarding the allocation of scarce resources and the interactions among these individuals and firms.
\ed

The goal of microeconomics is to analyze the market mechanisms that establish relative prices among goods and
services and allocate limited resources among alternative uses. Microeconomics shows conditions under which free
markets lead to desirable allocations. It also analyzes market failure, where markets fail to produce efficient
results. \v

Microeconomics focuses on supply and demand and other forces that determine price levels in the economy. It takes a
bottom-up approach to analyzing the economy. In other words, microeconomics tries to understand human choices,
decisions, and the allocation of resources. \v

Having said that, microeconomics does not try to answer or explain what forces should take place in a market. Rather,
it tries to explain what happens when there are changes in certain conditions. \v

In other words, microeconomics focuses on the study of individual markets, sectors, or industries as opposed to the
national economy as whole, which is studied in macroeconomics.

\bd[Macroeconomics]
\textbf{Macroeconomics} is a branch of economics dealing with performance,structure, behavior, and decision-making of an
economy as a whole.
\ed

Macroeconomics focuses on aggregates and econometric correlations, which is why governments and their agencies rely
on macroeconomics to formulate economic and fiscal policy. John Maynard Keynes is often credited as the founder of
macroeconomics, as he initiated the use of monetary aggregates to study broad phenomena. Some economists dispute his
theories, while many Keynesians disagree on how to interpret his work. \v

Microeconomics and macroeconomics are closely intertwined. Because changes in the overall economy arise from the
decisions of millions of individuals, it is impossible to understand macroeconomic developments without considering
the underlying microeconomic decisions. However, despite the inherent link between microeconomics and macroeconomics,
the two fields are distinct. Because they address different questions, each field has its own set of models. \v

A final, basic distinction of economics is between ``positive'' and ``normative'' economics.

\bd[Positive Economics]
\textbf{Positive economics} refers to the objective analysis in the study of economics.
\ed

Positive statements are descriptive. They make a claim about how the world is. Most economists look at what has
happened and what is currently happening in a given economy to form their basis of predictions for the future. This
investigative process is positive economics.

\bd[Normative Economics]
\textbf{Normative economics} is a perspective on economics that reflects normative, or ideologically prescriptive
judgments toward economic development, investment projects, statements, and scenarios.
\ed

Normative statements are prescriptive. They make a claim about how the world ought to be. Unlike positive economics,
which relies on objective data analysis, normative economics heavily concerns itself with value judgments and
statements of ``what ought to be'' rather than facts based on cause-and-effect statements. It expresses ideological
judgments about what may result in economic activity if public policy changes are made. Normative economic statements
can't be verified or tested. \v

A key difference between positive and normative statements is how we judge their validity. We can, in principle,
confirm or refute positive statements by examining evidence. By contrast, evaluating normative statements involves
values as well as facts. Deciding what is good or bad policy is not just a matter of science. It also involves our
views on ethics, religion, and political philosophy. \v

Positive and normative statements are fundamentally different, but within a person's set of beliefs, they are often
intertwined. In particular, positive views about how the world works affect normative views about what policies are
desirable. \v

Now, let's explore some more basic definitions of economics through introducing the fundamental principles of economics.

\section{The Ten Principles Of Economics}

The study of economics has many facets, but it is unified by several central ideas. In this chapter, we look at ``Ten
Principles of Economics''.

\subsection*{Principle 1: People Face Trade-offs}

To get something that we like, we usually have to give up something else that we also like. There is no such thing as
a free lunch (TINSTAAFL). Making decisions requires trading one goal for another.

\bd[Trade-off]
A \textbf{trade-off} is a situational decision that involves diminishing or losing one quality, quantity, or property
of a set or design in return for gains in other aspects.
\ed

In simple terms, a trade-off is where one thing increases, and another must decrease. A trade-off, then, involves a
sacrifice that must be made to obtain a certain product, service, or experience, rather than others that could be
made or obtained using the same required resources.

\be
Consider a student who must decide how to allocate her most valuable resource—her time. She can spend all of her time
studying economics, spend all of it studying psychology, or divide it between the two fields. For every hour she
studies one subject, she gives up an hour she could have used studying the other. And for every hour she spends
studying, she gives up an hour she could have spent napping, bike riding, playing video games, or working at her
part-time job for some extra spending money.
\ee

When people are grouped into societies, they face different kinds of trade-offs. \be One classic trade-off is between
``guns and butter''. The more a society spends on national defense (guns) to protect itself from foreign aggressors,
the less it can spend on consumer goods (butter) to raise its standard of living. \ee

A special example of a trade-off is the trade-off between (economic) efficiency and (economic) equality.

\bd[Economic Efficiency]
\textbf{Economic efficiency} is when all goods and factors of production in an economy are distributed or allocated
to their most valuable uses and waste is eliminated or minimized.
\ed

\bd[Economic Equality]
\textbf{Economic equality} is the property of distributing economic prosperity fairly among the members of society.
\ed

\bd[Equality-Efficiency Tradeoff]
An \textbf{equality-efficiency trade-off } is when there is some kind of conflict between maximizing economic
efficiency and maximizing the economic equality (or fairness) of society in some way.
\ed

When government policies are designed, these two goals often conflict. \v

Recognizing that people face trade-offs does not by itself tell us what decisions they will or should make.
Nonetheless, people are likely to make good decisions only if they understand the options available to them. Our
study of economics, therefore, starts by acknowledging life's trade-offs.

\subsection*{Principle 2: The Cost Of Something Is What You Give Up To Get It}

Because people face trade-offs, making decisions requires comparing the costs and benefits of alternative courses of
action. In many cases, however, the cost of an action is not as obvious as it might first appear.

\be
Consider the decision to go to college. The main benefits are intellectual enrichment and a lifetime of better job
opportunities. But what are the costs? To answer this question, you might be tempted to add up the money you spend on
tuition, books, room, and board. Yet this total does not truly represent what you give up to spend a year in college.
This calculation has two problems. First, it includes some things that are not really costs of going to college. Even
if you quit school, you need a place to sleep and food to eat. Room and board are costs of going to college only to
the extent that they exceed the cost of living and eating at home or in your own apartment. Second, this calculation
ignores the largest cost of going to college, your time. When you spend a year listening to lectures, reading
textbooks, and writing papers, you cannot spend that time working at a job and earning money. For most students, the
earnings they give up to attend school are the largest cost of their education.
\ee

In economics a trade-off is expressed in terms of the opportunity cost of a particular choice, which is the loss of
the most preferred alternative given up.

\bd[Opportunity Cost]
\textbf{Opportunity costs} represent the potential benefits an individual,investor, or business misses out on when
choosing one alternative over another.
\ed

The opportunity cost of an item is what you give up to get that item. When making any decision, decision makers
should take into account the opportunity costs of each possible action. \v

Because opportunity costs are, by definition, unseen, they can be easily overlooked. Understanding the potential
missed opportunities when a business or individual chooses one investment over another allows for better
decision-making.

\subsection*{Principle 3: Rational People Think At The Margin}

Economics focuses on the actions of human beings, based on assumptions that humans act with rational behavior,
seeking the most optimal level of benefit or utility.

\bd[Rational Behaviour]
\textbf{Rational behavior} refers to a decision-making process that is based on making choices that result in the
optimal level of benefit or utility for an individual.
\ed

The assumption of rational behavior implies that people would rather take actions that benefit them versus actions
that are neutral or harm them. Most classical economic theories are based on the assumption that all individuals
taking part in an activity are behaving rationally. Rational people systematically do the best they can to achieve
their objectives, given the available opportunities. \v

Rational people know that decisions in life are rarely black and white but often involve shades of gray. For example,
at dinnertime, you don't ask yourself ``Should I fast or eat like a pig?''. More likely, the question you face is
``Should I take that extra spoonful of mashed potatoes''? In economics this concept is being expressed through the
notion of ``marginal change''.

\bd[Marginal Change]
\textbf{Marginal change} is used to describe a small incremental adjustment to an existing plan of action.
\ed

Keep in mind that margin means ``edge'', so marginal changes are adjustments around the edges of what you are doing.
Rational people make decisions by comparing marginal benefits and costs.

\bd[Marginal Benefit]
A \textbf{marginal benefit} is a maximum amount a consumer is willing to pay for an additional good or service. It is
also the additional satisfaction or utility that a consumer receives when the additional good or service is purchased.
\ed

\bd[Marginal Cost]
A \textbf{marginal cost} is the small but measurable amount a producer has to pay for producing an additional unit of a
good or service.
\ed

The marginal benefit for a consumer tends to decrease as consumption of the good or service increases. If the
additional satisfaction obtained by an addition in the units of a good or a commodity is equal to the price a
consumer is willing to pay for that good he achieves maximum satisfaction, which is the main goal of every rational
consumer. A rational decision maker takes an action if and only if the action's marginal benefit exceeds its marginal
cost.

\be
For example, suppose you are considering watching a movie tonight. You pay \$40 a month for a movie streaming service
that gives you unlimited access to its film library, and you typically watch 8 movies a month. What cost should you
take into account when deciding whether to stream another movie? You might at first think the answer is \$40/8, or
\$5, which is the average cost of a movie. More relevant for your decision, however, is the marginal cost—the extra
cost that you would incur by streaming another film. Here, the marginal cost is zero because you pay the same \$40
for the service regardless of how many movies you stream. In other words, at the margin, streaming a movie is free.
The only cost of watching a movie tonight is the time it takes away from other activities.
\ee

Marginal decision making can explain some otherwise puzzling phenomena. Let's see on of the most famous ``
paradoxes'', the so called ``paradox of value'', also known as the ``diamond–water paradox''.

\be
Here is a classic question: Why is water so cheap, while diamonds are so expensive? Humans need water to survive,
while diamonds are unnecessary. Yet people are willing to pay much more for a diamond than for a cup of water. The
reason is that a person's willingness to pay for a good is based on the marginal benefit that an extra unit of the
good would yield. The marginal benefit, in turn, depends on how many units a person already has. Water is essential,
but the marginal benefit of an extra cup is small because water is plentiful. By contrast, no one needs diamonds to
survive, but because diamonds are so rare, the marginal benefit of an extra diamond is large.
\ee

\subsection*{Principle 4: People Respond To Incentives}

Because rational people make decisions by comparing costs and benefits, they respond to incentives.

\bd[Incentive]
\textbf{Incentive} is something that induces a person to act in a certain way. Incentives may possess a negative or a
positive intention.
\ed

Incentives play a central role in the study of economics, since they are key to analyzing how markets work.

\be
For example, when the price of apples rises, people decide to eat fewer apples. At the same time, apple orchards
decide to hire more workers and harvest more apples. In other words, a higher price in a market provides an incentive
for buyers to consume less and an incentive for sellers to produce more.
\ee

Public policymakers should never forget about incentives: Many policies change the costs or benefits that people face
and, as a result, alter their behavior. When policymakers fail to consider how their policies affect incentives,they
often face unintended consequences. \v

The first four principles discussed how individuals make decisions. As we go about our lives, many of our decisions
affect not only ourselves but other people as well. The next three principles concern how people interact with one
another.

\subsection*{Principle 5: Trade Can Make Everyone Better Off}

There are two ways to compare the abilities of two people to produce a good. The first one is the so called
``absolute advantage''.

\bd[Absolute Advantage]
\textbf{Absolute advantage} is the ability of an individual, company, region, or country to produce a greater
quantity of a good or service with the same quantity of inputs per unit of time, or to produce the same quantity of a
good or service per unit of time using a lesser quantity of inputs, than its competitors.
\ed

The person who can produce the good with the smaller quantity of inputs is said to have an absolute advantage in
producing the good. Absolute advantage can be accomplished by creating the good or service at a lower absolute cost
per unit using a smaller number of inputs, or by a more efficient process.

\bd[Unit Cost]
A \textbf{unit cost} is a total expenditure incurred by a company to produce,store,and sell one unit of a particular
product or service.
\ed

The second one is the so called ``comparative advantage''.

\bd[Comparative Advantage]
\textbf{Comparative advantage} is an economy's ability to produce a particular good or service at a lower opportunity
cost than its trading partners.
\ed

The person who has the lower opportunity cost of producing the good is said to have a comparative advantage.
Comparative advantage is used to explain why companies, countries, or individuals can benefit from trade, since the
gains from trade are based on comparative advantage, not absolute advantage.

\bd[Trade]
\textbf{Trade} is a basic economic concept involving the buying and selling of goods and services, with compensation
paid by a buyer to a seller, or the exchange of goods or services between parties.
\ed

Trade can take place within an economy between producers and consumers. However, it can also take place between
countries in international trade. International trade allows countries to expand markets for both goods and services
that otherwise may not have been available. As a result the market contains greater competition and therefore, more
competitive prices, which brings a cheaper product home to the consumer. \v

In general trade makes everyone better off because it allows people to specialize in those activities in which they
have a comparative advantage. In a broader scope, trade allows for specialization in products that benefits countries.
Interdependence and trade are desirable because they allow everyone to enjoy a greater quantity and variety of
goods and services. For centuries, it was widely believed that in international trade one country's gain from an
exchange must be the other country's loss. However, trade is not like a sports competition, where one side gains and
the other side loses.

\subsection*{Principle 6: Markets Are Usually A Good Way To Organize Economic Activity}

\bd[Social System]
A \textbf{social system} is the patterned network of relationships constituting a coherent whole that exist between
individuals, groups, and institutions. It is the formal structure of role and status that can form in a small, stable
group.
\ed

\be
Examples of social systems include nuclear family units, communities,cities, nations, college campuses, corporations,
and industries.
\ee

An economic system is a type of social system which confront and solve the four fundamental economic problems:
\bit
\item What kinds and quantities of goods shall be produced.
\item How goods shall be produced.
\item How the output will be distributed.
\item When to produce.
\eit

\bd[Economic System]
An \textbf{economic system} is a system of production, resource allocation and distribution of goods and services
within a society or a given geographic area. It includes the combination of the various institutions, agencies,
entities, decision-making processes and patterns of consumption that comprise the economic structure of a given
community.
\ed

There are many different economic systems. In what follows we will mention the most important ones.

\bd[Centrally Planned Economy]
A \textbf{centrally planned economy}, also known as a command economy, is an economic system in which a central
authority, such as a government, makes economic decisions regarding the manufacturing and the distribution of
products.
\ed

The production of goods and services in centrally planned economy is often done by state-owned enterprises, which are
government owned companies. In centrally planned economies prices are controlled by bureaucrats. \v

Centrally planned economies have failed because they did not allow the market to work. Most countries that once had
centrally planned economies have abandoned the system and instead have adopted market economies in which such
decisions are traditionally made by businesses and consumers.

\bd[Market]
A\textbf{market} is a place where parties can gather to facilitate the exchange of goods and services. The parties
involved are usually buyers (consumers) and sellers (producers).
\ed

A market may be physical like a retail outlet, where people meet face-to-face, or virtual like an online market,
where there is no direct physical contact between buyers and sellers.

\bd[Market Economy]
A \textbf{market economy} is an economic system in which the decisions regarding investment, production and
distribution are guided by the price signals created by the forces of the market.
\ed

In a market economy, the decisions of a central planner are replaced by the decisions of millions of firms and
households. Firms decide whom to hire and what to make. Households decide which firms to work for and what to buy
with their incomes. These firms and households interact in the marketplace, where prices and self-interest guide
their decisions. Because a market economy rewards people for their ability to produce things that other people are
willing to pay for, there will be an unequal distribution of economic prosperity. \v

Market economies range from minimally regulated free market systems where state activity is restricted to providing
public goods and services and safeguarding private ownership, to interventionist forms where the government plays an
active role in serving special interests and promoting social welfare.

\bd[Free Market Economy]
A \textbf{free market economy} is an economic system in which economic decisions and the pricing of goods and
services are guided by the interactions of a country's individual citizens and businesses. There may be some
government intervention or central planning, but usually this term refers to an economy that is more market oriented
in general.
\ed

At first glance, the success of market economies is puzzling. In a market economy, no one is looking out for the
well-being of society as a whole. Free markets contain many buyers and sellers of numerous goods and services, and
all of them are interested primarily in their own well-being. Yet despite decentralized decision making and
self-interested decision makers, market economies have proven remarkably successful in organizing economic activity
to promote overall prosperity. Adam Smith's 1776 work suggested that although individuals are motivated by
self-interest, an invisible hand guides this self-interest into promoting society's economic well-being.

\bd[Invisible Hand]
The \textbf{invisible hand} is a metaphor for the unseen forces that move the free market economy. Through individual
self-interest and freedom of production and consumption, the best interest of society, as a whole, are fulfilled.
\ed

Prices are the instrument with which the invisible hand directs economic activity.

\bd[Market Price]
The \textbf{market price} is the current price at which an asset or service can be bought or sold.
\ed

The market price of an asset or service is determined by the forces of supply and demand. The price at which quantity
supplied equals quantity demanded is the market price. In any market, buyers look at the price when deciding how much
to demand, and sellers look at the price when deciding how much to supply. As a result of these decisions, market
prices reflect both the value of a good to society and the cost to society of making the good. \v

Smith's great insight was that prices adjust to guide buyers and sellers to reach outcomes that, in many cases,
maximize the well-being of society as a whole. Smith's insight has an important corollary: When a government prevents
prices from adjusting naturally to supply and demand, it impedes the invisible hand's ability to coordinate the
decisions of the households and firms that make up an economy. \v

Finally, there is another important economic system called ``capitalism''.

\bd[Capitalism]
\textbf{Capitalism} is an economic system based on the private ownership and control of the means of production and
their operation for profit. Central characteristics of capitalism include capital accumulation,competitive markets, a
price system determined by supply and demand, private property, property rights recognition, voluntary exchange, and
wage labor
\ed

In a capitalist market economy, decision-making and investments are determined by owners of wealth, property, ability
to maneuver capital or production ability in capital and financial markets—whereas prices and the distribution of
goods and services are mainly determined by competition in markets.

\subsection*{Principle 7: Governments Can Sometimes Improve Market Outcomes}

There are two broad reasons for the government to interfere with the economy. First reason we need government is that
the invisible hand can work its magic only if the government enforces the rules and maintains the institutions that
are key to a market economy. Most important, market economies need institutions to enforce property rights so
individuals can own and control scarce resources.

\bd[Property Rights]
\textbf{Property rights} define the theoretical and legal ownership of resources and how they can be used. These
resources can be both tangible or intangible and can be owned by individuals, businesses, and governments.
\ed

We all rely on government, provided police and courts to enforce our rights over the things we produce, and the
invisible hand counts on our ability to enforce those rights. \v

The second reason we need government is that, although the invisible hand is powerful, it is not omnipotent. There
are two broad rationales for a government to intervene in the economy and change the allocation of resources that
people would choose on their own: to promote efficiency or to promote equality. That is, most policies aim either to
enlarge the economic pie or to change how the pie is divided. \v

Consider first the goal of efficiency. Although the invisible hand usually leads markets to allocate resources to
maximize the size of the economic pie, this is not always the case. Economists use the term market failure to refer
to a situation in which the market on its own fails to produce an efficient allocation of resources.

\bd[Market Failure]
\textbf{Market failure} is a situation defined by an inefficient distribution of goods and services in the free
market. In market failure, the individual incentives for rational behavior do not lead to rational outcomes for the
group.
\ed

In other words, in market failure, each individual makes the correct decision for themselves, but those prove to be
the wrong decisions for the group. In traditional microeconomics, this can sometimes be shown as a steady-state
disequilibrium in which the quantity supplied does not equal the quantity demanded. \v

As we will see, one possible cause of market failure is an externality, which is the impact of one person's actions
on the well-being of a bystander.

\bd[Externality]
An \textbf{externality} is a cost or benefit caused by a producer that is not financially incurred or received by that
producer.
\ed

An externality can be both positive or negative and can stem from either the production or consumption of a good or
service. The costs and benefits can be both private—to an individual or an organization—or social, meaning it can
affect society as a whole. Externalities by nature are generally environmental, such as natural resources or public
health.

\be
The classic example of an externality is pollution. When the production of a good pollutes the air and creates health
problems for those who live near the factories, the market on its own may fail to take this cost into account.
\ee

Another possible cause of market failure is market power, which refers to the ability of a single person or firm (or
a small group of them) to unduly influence market prices.

\bd[Market Power]
\textbf{Market power} refers to a single person's or company's relative ability to manipulate the price of an item in
the marketplace by manipulating the level of supply, demand or both.
\ed

A company with substantial market power has the ability to manipulate the market price and thereby control its profit
margin, and possibly the ability to increase obstacles to potential new entrants into the market. Firms that have
market power are often described as ``price makers'' because they can establish or adjust the marketplace price of an
item without relinquishing market share.

\be
If everyone in town needs water but there is only one well, the owner of the well does not face the rigorous
competition with which the invisible hand normally keeps self-interest in check.
\ee

In the presence of externalities or market power, well-designed public policy can enhance economic efficiency. \v

Now consider the goal of equality. Even when the invisible hand yields efficient outcomes, it can nonetheless leave
sizable disparities in economic well-being. A market economy rewards people according to their ability to produce
things that other people are willing to pay for. The world's best basketball player earns more than the world's best
chess player simply because people are willing to pay more to watch basketball than chess. The invisible hand does
not ensure that everyone has sufficient food, decent clothing, and adequate healthcare. This inequality may,
depending on one's political philosophy, call for government intervention. In practice, many public policies, such as
the income tax and the welfare system, aim to achieve a more equal distribution of economic well-being. \v

To say that the government can improve market outcomes does not mean that it always will. Public policy is made not
by angels but by a political process that is far from perfect. Sometimes policies are designed to reward the
politically powerful. Sometimes they are made by well-intentioned leaders who are not fully informed. As you study
economics, you will become a better judge of when a government policy is justifiable because it promotes efficiency
or equality and when it is not. \v

We started by discussing how individuals make decisions and then looked at how people interact with one another. All
these decisions and interactions together make up ``the economy''. The last three principles concern the workings of
the economy as a whole.

\subsection*{Principle 8: A Country's Standard Of Living Depends On Its Ability To Produce Goods And Services}

Differences in the standard of living from one country to another are quite large. Changes in living standards over
time are also quite large. The explanation for differences in living standards lies in differences in productivity.

\bd[Productivity]
\textbf{Productivity} is the quantity of goods and services produced by each unit of labor input.
\ed

Measurements of productivity are often expressed as a ratio of an aggregate output to a single input or an aggregate
input used in a production process. \v

High productivity implies a high standard of living. The relationship between productivity and living standards is
simple, but its implications are far-reaching. If productivity is the primary determinant of living standards, other
explanations must be less important. The relationship between productivity and living standards also has profound
implications for public policy. When thinking about how any policy will affect living standards, the key question is
how it will affect our ability to produce goods and services. To boost living standards, policymakers need to raise
productivity by ensuring that workers are well educated, have the tools they need to produce goods and services, and
have access to the best available technology.

\subsection*{Principle 9: Prices Rise When The Government Prints Too Much Money}

\bd[Purchasing Power]
\textbf{Purchasing power} is the value of a currency expressed in terms of the number of goods or services that one unit
of money can buy.
\ed

Purchasing power is important because, all else being equal, inflation decreases the number of goods or services you
would be able to purchase.

\bd[Inflation]
\textbf{Inflation} is the decline of purchasing power of a given currency over time.
\ed

A quantitative estimate of the rate at which the decline in purchasing power occurs can be reflected in the increase
of an average price level of a basket of selected goods and services in an economy over some period of time. The rise
in the general level of prices, often expressed as a percentage, means that a unit of currency effectively buys less
than it did in prior periods. \v

What causes inflation? In almost all cases of large or persistent inflation, the culprit is growth in the quantity of
money. When a government creates large quantities of the nation's money, the value of the money falls, because the
individuals have more money and the demand increases. When the demand increases, price also increases and creates
inflation of money.

\subsection*{Principle 10: Society Faces A Tradeoff Between Inflation And Unemployment}

\bd[Monetary Authority / Central Bank]
A \textbf{monetary\footnote{Monetary means ``relating to money or currency''.}authority} or \textbf{central bank} is
the entity that manages a country's currency and money supply, often with the objective of controlling inflation,
interest rates, real GDP or unemployment rate.
\ed

With its monetary tools, a monetary authority is able to effectively influence parameters which control the cost and
availability of money (like interest rates). \v

In contrast to a commercial bank, a central bank possesses a monopoly on increasing the monetary base.

\bd[Financial Institution]
A \textbf{financial institution} is a company engaged in the business of dealing with financial and monetary
transactions such as deposits, loans, investments, and currency exchange.
\ed

Financial institutions encompass a broad range of business operations within the financial services sector including
banks, trust companies, insurance companies, brokerage firms, and investment dealers. Virtually everyone living in a
developed economy has an ongoing or at least periodic need for the services of financial institutions.

\bd[Commercial Bank]
A \textbf{commercial bank} is a financial institution which accepts deposits from the public and gives loans for the
purposes of consumption and investment to make profit.
\ed

\bd[Monetary Policy]
\textbf{Monetary policy} is the policy adopted by the monetary authority of a nation to control either the interest
rate payable for very short-term borrowing (borrowing by banks from each other to meet their short-term needs) or the
money supply, often as an attempt to reduce inflation or the interest rate, to ensure price stability and general
trust of the value and stability of the nation's currency.
\ed

The goal of a monetary policy is to keep the economy humming along at a rate that is neither too hot nor too cold.
The central bank may force up interest rates on borrowing in order to discourage spending or force down interest
rates to inspire more borrowing and spending. The main weapon at its disposal is the nation's money. The central bank
sets the rates it charges to loan money to the nation's banks. When it raises or lowers its rates, all financial
institutions tweak the rates they charge all of their customers, from big businesses borrowing for major projects to
home buyers applying for mortgages. \v

While an increase in the quantity of money primarily raises prices in the long run, the short-run story is more
complex. In short run increasing the amount of money in the economy stimulates the overall level of spending and thus
the demand for goods and services. Higher demand may over time cause firms to raise their prices, but in the
meantime, it also encourages them to hire more workers and produce a larger quantity of goods and services. Finally,
more hiring means lower unemployment. This line of reasoning leads to one final economy-wide trade-off: a short-run
trade-off between inflation and unemployment. \v

Although some economists still question these ideas, most accept that society faces a short-run trade-off between
inflation and unemployment. This simply means that, over a period of a year or two, many economic policies push
inflation and unemployment in opposite directions. Policymakers face this trade-off regardless of whether inflation
and unemployment both start out at high levels (as they did in the early 1980s), at low levels (as they did in the
late 1990s), or someplace in between. This short-run trade-off plays a key role in the analysis of the business
cycle—the irregular and largely unpredictable fluctuations in economic activity, as measured by the production of
goods and services or the number of people employed.

\bd[Business Cycles]
\textbf{Business cycles} are intervals of expansion followed by recession in economic activity.
\ed

Policymakers can exploit the short-run trade-off between inflation and unemployment using various policy instruments.
By changing the amount that the government spends, the amount it taxes, and the amount of money it prints,
policymakers can influence the overall demand for goods and services. Changes in demand in turn influence the
combination of inflation and unemployment that the economy experiences in the short run. Because these instruments of
economic policy are so powerful, how policymakers should use them to control the economy, if at all, is a subject of
continuing debate.